\chapter{The console}\label{ch:console}

The page has a second face. The tabs are made for a desk: a mouse, a large
window, time to read. The console is made for the other case --- a small screen,
gloves, noise, a person who has to decide now. It is opened from the status bar
and it covers the page completely, because a half-covered page is a page that can
be misread.

Two rules shaped it. The first is that the console must never be a dead end: a
view that has nothing to report still says what it has, because an operator under
pressure reads ``nothing outstanding'' as ``the tool has stopped working''. The
second is that the console carries the same data as the page, not a summary of it:
every figure in it comes from the same function the tab calls.

\section{What the console shows}
Eight views, one per area of the page. Each view builds its own items from the
loaded data; none of them filters a shared list, which is why two views never show
the same thing twice. The counts below are from the corpus used throughout this
part.

{\small
\begin{longtable}{@{}p{3.0cm}p{2.1cm}p{6.8cm}p{3.8cm}@{}}
\caption{The eight console views and what each built from the corpus.}\label{tab:console-views} \\
\toprule \textbf{View} & \textbf{Items} & \textbf{Kinds of item} & \textbf{Severity} \\\midrule\endfirsthead
\toprule \textbf{View} & \textbf{Items} & \textbf{Kinds of item} & \textbf{Severity} \\\midrule\endhead
\texttt{load} & 4 & file 3, summary 1 & ok 4 \\
\texttt{sop} & 20 & rule 19, summary 1 & none 19, ok 1 \\
\texttt{results} & 4 & results 3, summary 1 & ok 4 \\
\texttt{panel} & 54 & chart 53, summary 1 & none 53, ok 1 \\
\texttt{graphs} & 26 & graph 25, summary 1 & none 25, ok 1 \\
\texttt{review} & 13 & finding 12, summary 1 & watch 10, alarm 3 \\
\texttt{export} & 12 & export 12 & ok 12 \\
\texttt{integrity} & 1 & integrity 1 & ok 1 \\
\bottomrule
\end{longtable}}


\section{An item}
An item is one fact with everything needed to act on it: a code, a title, a value
with its unit, the limits it was judged against, a sentence of state, the evidence
behind it, a small table of detail and the actions that can follow. The item below
is the first one the load view builds.

{\small
\begin{longtable}{@{}ll@{}}
\toprule \textbf{Field} & \textbf{Content} \\\midrule\endfirsthead
\toprule \textbf{Field} & \textbf{Content} \\\midrule\endhead
\texttt{code} & LOAD \\
\texttt{title} & Files read into this session \\
\texttt{value} & 3 files \\
\texttt{kind} & files read \\
\texttt{state} & Decoded file metadata and stored results are ready. Open File integrity for any … \\
\bottomrule
\end{longtable}}

\vhead{The detail rows of that item}
{\small
\begin{longtable}{@{}ll@{}}
\toprule \textbf{Row} & \textbf{Value} \\\midrule\endfirsthead
\toprule \textbf{Row} & \textbf{Value} \\\midrule\endhead
Stored results & 830 \\
Curves decoded & 751 \\
Melting results & 79 \\
Instruments & SVT001, Pilot-510, DVT07 \\
Operators & Demo \\
Active profile & SC screening — SOP-06 laboratory default v1.0 \\
\bottomrule
\end{longtable}}


\section{Pictograms}
A console is read at a glance, so every item carries a mark for its kind. The
marks are drawn from the characters the font already has, which keeps the page a
single file with no image to load.

{\small\noindent\makebox[\textwidth][l]{\begin{tabular}{@{}ll@{}}
\toprule \textbf{Kind} & \textbf{Mark} \\\midrule
\texttt{file} & ▤ \\
\texttt{summary} & ▣ \\
\texttt{rule} & § \\
\texttt{sample} & ◉ \\
\texttt{run} & ▥ \\
\texttt{chart} & ◪ \\
\texttt{graph} & ◩ \\
\texttt{finding} & ⚠ \\
\texttt{export} & ↓ \\
\texttt{baseline} & ◐ \\
\texttt{instrument} & ⚙ \\
\texttt{control} & ◈ \\
\bottomrule
\end{tabular}}}


\section{Working without a mouse}
Everything in the console can be reached by typing. The recogniser list below is
the one the code carries; a button in the interface sends exactly the same text,
so the two paths cannot drift apart.

{\small
\begin{longtable}{@{}ll@{}}
\caption{Every command the console recognises.}\label{tab:console-commands} \\
\toprule \textbf{Command} & \textbf{Pattern} \\\midrule\endfirsthead
\toprule \textbf{Command} & \textbf{Pattern} \\\midrule\endhead
next & \texttt{\textasciicircum{}next\$} \\
previous & \texttt{\textasciicircum{}previous\$} \\
chart previous & \texttt{\textasciicircum{}chart previous\$} \\
chart next & \texttt{\textasciicircum{}chart next\$} \\
chart close & \texttt{\textasciicircum{}chart close\$} \\
zoom & \texttt{\textasciicircum{}zoom (in|out|reset)\$} \\
pan & \texttt{\textasciicircum{}pan (left|right|up|down)\$} \\
chart type menu & \texttt{\textasciicircum{}(chart type|type menu|change chart type)\$} \\
x axis menu & \texttt{\textasciicircum{}(x axis|axis menu|change x axis)\$} \\
close menu & \texttt{\textasciicircum{}(close|close menu|dismiss)\$} \\
read staged files & \texttt{\textasciicircum{}read files\$} \\
choose experiment files & \texttt{\textasciicircum{}choose files\$} \\
leave for a page & \texttt{\textasciicircum{}open page (load|sop|results|panel|graphs|review|export…} \\
console view & \texttt{\textasciicircum{}(view )?(load files|load view|sop view|results view|pa…} \\
chart type & \texttt{\textasciicircum{}type (i|ewma|cusum|trend|xbars|p|u|c|funnel)\$} \\
x axis & \texttt{\textasciicircum{}x (order|date|hours|cycles)\$} \\
theme light & \texttt{\textasciicircum{}theme light\$} \\
theme dark & \texttt{\textasciicircum{}theme dark\$} \\
chart & \texttt{\textasciicircum{}chart\$} \\
panel & \texttt{\textasciicircum{}panel\$} \\
results & \texttt{\textasciicircum{}results\$} \\
review & \texttt{\textasciicircum{}review\$} \\
load & \texttt{\textasciicircum{}load\$} \\
sop & \texttt{\textasciicircum{}sop\$} \\
next run & \texttt{\textasciicircum{}next run\$} \\
download a file & \texttt{\textasciicircum{}export (\textbackslash\{\}d+)\$} \\
exit & \texttt{\textasciicircum{}exit\$} \\
\bottomrule
\end{longtable}}


\section{The actions offered on the items}
An item that can be acted on carries its action with it, so the operator never has
to remember where to go next.

{\small
\begin{longtable}{@{}ll@{}}
\toprule \textbf{Offered as} & \textbf{Sends} \\\midrule\endfirsthead
\toprule \textbf{Offered as} & \textbf{Sends} \\\midrule\endhead
Choose experiment files & \texttt{choose files} \\
Open the SOP view on the page & \texttt{open page sop} \\
Open the results view on the page & \texttt{open page results} \\
Open the control panel on the page & \texttt{open page panel} \\
Show the chart & \texttt{chart} \\
Next run & \texttt{next run} \\
Open the graphs view on the page & \texttt{open page graphs} \\
Draw the graph & \texttt{chart} \\
Open the review view on the page & \texttt{open page review} \\
Download Cq values (CSV) & \texttt{export 0} \\
Open the formats view on the page & \texttt{open page export} \\
Download this file & \texttt{export 0} \\
Download this file & \texttt{export 1} \\
Download this file & \texttt{export 2} \\
Download this file & \texttt{export 3} \\
Download this file & \texttt{export 4} \\
Download this file & \texttt{export 5} \\
Download this file & \texttt{export 6} \\
Download this file & \texttt{export 7} \\
Download this file & \texttt{export 8} \\
Download this file & \texttt{export 9} \\
Download this file & \texttt{export 10} \\
Open File integrity on the page & \texttt{open page integrity} \\
\bottomrule
\end{longtable}}


\section{The view builders}
Each view is a function of the loaded data and nothing else. Two are printed here;
the rest are in Chapter~\ref{ch:val-functions} with the record of what they
returned when they were called.

\vhead{\vfn{mgViewLoad}}
\begin{lstlisting}[language=JavaScript,basicstyle=\ttfamily\scriptsize\color{white},numbers=left,firstnumber=7449]
function mgViewLoad(){
  const out=[];
  if(!RUNS.length&&!PASTED){
    const staged=Array.isArray(STAGED)&&STAGED.length;
    return [mgIt({kind:"intake",sev:"none",pict:MG_PICT.file,kindWord:"intake",code:"LOAD",
      title:staged?"Files staged, ready to read":"No experiment files are loaded",
      value:staged?String(STAGED.length):"0",unit:staged?"files staged":"files",
      state:staged?"The files passed intake. Read them to decode runs and build the views.":"Drop .ixo, .eds, .edt or .zip files on the load view to begin.",
      actions:staged
        ?[{key:"3",label:"Read staged files",cmd:"read files"},{key:"4",label:"Choose another selection",cmd:"choose files"}]
        :[{key:"3",label:"Choose experiment files",cmd:"choose files"}]})];
  }
  const meta=mgSafe(()=>metaRows(),[]);
  if(RUNS.length){
    const results=RUNS.reduce((a,r)=>a+((r.wells||[]).length),0);
    const curves=RUNS.reduce((a,r)=>a+mgSafe(()=>uniqueCurveCount(r),0),0);
    const tm=RUNS.reduce((a,r)=>a+((r.tmWells||[]).length),0);
    /* A file that could not be read, or whose container checksum does not match, is what
       raises this view. Without the test the view threw and the operator saw nothing. */
    const bad=RUNS.some(r=>r.acqError||(r.integrity&&r.integrity.ok===false));
    out.push(mgIt({kind:"summary",sev:bad?"alarm":"ok",pict:MG_PICT.summary,kindWord:"files read",code:"LOAD",
      title:"Files read into this session",value:String(RUNS.length),unit:RUNS.length===1?"file":"files",
      limits:`${results.toLocaleString()} stored results · ${curves.toLocaleString()} curves decoded`,
      state:"Decoded file metadata and stored results are ready. Open File integrity for any file-integrity mismatch.",
      evidence:uniq(RUNS.map(r=>r.platform||"LightCycler 480")).join(" · "),
      rows:[["Stored results",results.toLocaleString()],["Curves decoded",curves.toLocaleString()],
            ["Melting results",tm.toLocaleString()],
            ["Instruments",uniq(RUNS.map(r=>(r.meta&&r.meta.InstrumentName)||"not named")).join(", ")],
            ["Operators",uniq(meta.map(m=>m.operator).filter(Boolean)).join(", ")||"—"],
            ["Active profile",`${SOP.name} v${SOP.version}`]],
      actions:[{key:"3",label:"Choose experiment files",cmd:"choose files"}]}));
    RUNS.forEach((r,i)=>{
      const m=meta[i]||{},ok=true;
      out.push(mgIt({kind:"file",sev:r.acqError?"watch":"ok",pict:MG_PICT.file,kindWord:"file",
        code:"F"+(i+1),title:m.experiment||r.file||("Experiment "+(i+1)),
        value:String(m.quantification_results||0),unit:"stored results",
        limits:`${mgTxt(m.plate)} plate · ${mgTxt(m.cycles)} cycles · ${m.curves_decoded||0} curves`,
        state:r.acqError?("Fluorescence could not be read: "+r.acqError)
              :r.isTemplate?"Template or unrun experiment — plate setup, protocol and analysis settings only."
              :"File decoded; stored values are ready for review.",
        evidence:`${mgTxt(m.platform)} · ${m.instrument||"instrument not named"} · ${m.run_started||m.experiment_created||"no date"}`,
        rows:[["File",mgTxt(r.file)],["Instrument",mgTxt(m.instrument)],["Software",mgTxt(m.software_version)],
              ["Operator",mgTxt(m.operator)],["Run started",mgTxt(m.run_started)],["Run ended",mgTxt(m.run_ended)],
              ["Run minutes",mgTxt(m.run_duration_minutes)],["Channels",mgTxt(m.channels)],
              ["Analyses",String(m.analyses||0)],["Melting results",String(m.melting_results||0)],
              ["Bytes",m.source_bytes?Number(m.source_bytes).toLocaleString():"—"]],
        actions:[{key:"3",label:"Choose experiment files",cmd:"choose files"}]}));
    });
  }
  if(PASTED)out.push(mgIt({kind:"file",sev:"ok",pict:MG_PICT.file,kindWord:"pasted table",code:"CSV",
    title:"Pasted Cq table",value:String((PASTED.rows||[]).length),unit:"Cq values",
    limits:"no fluorescence curves in a pasted table",
    state:"Replicate checks and the inhibition test use these values; curve-based exports need an .ixo or .eds file.",
    evidence:PASTED.source||"pasted data"}));
  return out;
}
\end{lstlisting}

\vhead{\vfn{mgViewReview}}
\begin{lstlisting}[language=JavaScript,basicstyle=\ttfamily\scriptsize\color{white},numbers=left,firstnumber=7756]
function mgViewReview(){
  const events=mgSafe(()=>reviewEventsWithoutIntegrity()||[],[]);
  if(!events.length)return [mgIt({kind:"finding",sev:"ok",pict:MG_PICT.finding,kindWord:"review",code:"QC",
    title:"No finding was raised",value:"0",unit:"findings",
    limits:`${RUNS.length} run(s) examined`,
    state:"The file, result, timeline and interpretation checks all passed on the loaded runs.",
    actions:[{key:"3",label:"Open the review view on the page",cmd:"open page review"}]})];
  const err=events.filter(e=>e.severity==="error").length,rev=events.filter(e=>e.severity==="review").length;
  const by=new Map();events.forEach(e=>{const k=e.area||"other";const l=by.get(k)||[];l.push(e);by.set(k,l);});
  const out=[mgIt({kind:"summary",sev:err?"alarm":rev?"watch":"ok",pict:MG_PICT.summary,kindWord:"review",code:"QC",
    title:"Quality and forensic review",value:String(events.length),unit:"findings",
    limits:`${err} error · ${rev} review · ${by.size} area(s)`,
    state:err?`${err} finding(s) are error level: the evidence contradicts what the file claims about itself.`
      :rev?`${rev} finding(s) are for review; none is error level.`:"Every check passed.",
    evidence:[...by.keys()].slice(0,5).join(" · "),
    rows:[...by].map(([a,l])=>[a,`${l.filter(e=>e.severity==="error").length} error · ${l.filter(e=>e.severity==="review").length} review`]),
    actions:[{key:"3",label:"Open the review view on the page",cmd:"open page review"}]})];
  const rank=l=>l.some(e=>e.severity==="error")?0:1;
  [...by].sort((a,b)=>rank(a[1])-rank(b[1])).forEach(([area,list])=>{
    const e0=list.find(e=>e.severity==="error")||list[0];
    const errs=list.filter(e=>e.severity==="error").length;
    out.push(mgIt({kind:"finding",sev:errs?"alarm":"watch",pict:MG_PICT.finding,kindWord:"finding",
      code:errs?"ERR":"REV",title:area,value:String(list.length),unit:list.length===1?"finding":"findings",
      limits:`${errs} error · ${list.length-errs} review · expected: none`,
      state:e0.finding||"",evidence:e0.experiment||"",
      rows:list.slice(0,14).map(e=>[String(e.experiment||e.finding||"").slice(0,44),
        String(e.evidence||e.finding||"").slice(0,80)]),
      actions:[{key:"3",label:"Open the review view on the page",cmd:"open page review"}]}));
  });
  return out;
}
\end{lstlisting}

